\documentclass[UTF8,12pt]{ctexart}
\usepackage[paperwidth=240mm,paperheight=200mm,margin=14mm,top=8mm]{geometry}
\usepackage{amsmath,amssymb,mathtools}
\usepackage{xcolor}
\pagestyle{empty}
\setlength{\parindent}{0pt}
\setlength{\parskip}{0pt}
\newcommand{\qn}[1]{\textbf{\large #1}\ }
\xeCJKDeclareCharClass{CJK}{"2460 -> "24FF}
\begin{document}
\qn{2020.1}
当 $x \to 0^+$ 时，下列无穷小量中最高阶的是（　　）．
\\
\noindent (A) $\displaystyle\int_{0}^{x} (\mathrm{e}^{t^2} - 1)\,\mathrm{d}t$
\\
\noindent (B) $\displaystyle\int_{0}^{x} \ln(1 + \sqrt{t^3})\,\mathrm{d}t$
\\
\noindent (C) $\displaystyle\int_{0}^{\sin x} \sin t^2\,\mathrm{d}t$
\\
\noindent (D) $\displaystyle\int_{0}^{1-\cos x} \sqrt{\sin^3 t}\,\mathrm{d}t$

\vfill
\newpage
\qn{2020.2}
函数 $f(x) = \dfrac{\mathrm{e}^{\frac{1}{x-1}}\ln|1+x|}{(\mathrm{e}^x - 1)(x-2)}$ 的第二类间断点的个数为（　　）．
\\
\noindent (A) 1
\\
\noindent (B) 2
\\
\noindent (C) 3
\\
\noindent (D) 4

\vfill
\newpage
\qn{2020.3}
$\displaystyle\int_{0}^{1} \dfrac{\arcsin\sqrt{x}}{\sqrt{x(1-x)}}\,\mathrm{d}x =$（　　）．
\\
\noindent (A) $\dfrac{\pi^2}{4}$
\\
\noindent (B) $\dfrac{\pi^2}{8}$
\\
\noindent (C) $\dfrac{\pi}{4}$
\\
\noindent (D) $\dfrac{\pi}{8}$

\vfill
\newpage
\qn{2020.4}
设 $f(x) = x^2\ln(1-x)$，当 $n \geqslant 3$ 时，$f^{(n)}(0) =$（　　）．
\\
\noindent (A) $-\dfrac{n!}{n-2}$
\\
\noindent (B) $\dfrac{n!}{n-2}$
\\
\noindent (C) $-\dfrac{(n-2)!}{n}$
\\
\noindent (D) $\dfrac{(n-2)!}{n}$

\vfill
\newpage
\qn{2020.5}
关于函数 $f(x,y) = \begin{cases} xy, & xy \neq 0, \\ x, & y = 0, \\ y, & x = 0, \end{cases}$ 给出如下结论：
① $\left.\dfrac{\partial f}{\partial x}\right|_{(0,0)} = 1$；
② $\left.\dfrac{\partial^2 f}{\partial x \partial y}\right|_{(0,0)} = 1$；
③ $\lim_{(x,y)\to(0,0)} f(x,y) = 0$；
④ $\lim_{y\to 0}\lim_{x\to 0} f(x,y) = 0$．
其中正确的个数是（　　）．
\\
\noindent (A) 4
\\
\noindent (B) 3
\\
\noindent (C) 2
\\
\noindent (D) 1

\vfill
\newpage
\qn{2020.6}
设函数 $f(x)$ 在 $[-2,2]$ 上可导，且 $f'(x) > f(x) > 0$，则（　　）．
\\
\noindent (A) $\dfrac{f(-2)}{f(-1)} > 1$
\\
\noindent (B) $\dfrac{f(0)}{f(-1)} > \mathrm{e}$
\\
\noindent (C) $\dfrac{f(1)}{f(-1)} < \mathrm{e}^2$
\\
\noindent (D) $\dfrac{f(2)}{f(-1)} < \mathrm{e}^3$

\vfill
\newpage
\qn{2020.7}
设 4 阶矩阵 $\boldsymbol{A} = (a_{ij})_{4\times 4}$ 不可逆，$a_{12}$ 的代数余子式 $A_{12} \neq 0$，$\boldsymbol{\alpha}_1, \boldsymbol{\alpha}_2, \boldsymbol{\alpha}_3, \boldsymbol{\alpha}_4$ 为矩阵 $\boldsymbol{A}$ 的列向量组，$\boldsymbol{A}^*$ 为 $\boldsymbol{A}$ 的伴随矩阵，则方程组 $\boldsymbol{A}^*\boldsymbol{X} = \boldsymbol{0}$ 的通解为（　　）．
\\
\noindent (A) $\boldsymbol{X} = k_1\boldsymbol{\alpha}_1 + k_2\boldsymbol{\alpha}_2 + k_3\boldsymbol{\alpha}_3$，其中 $k_1, k_2, k_3$ 为任意常数
\\
\noindent (B) $\boldsymbol{X} = k_1\boldsymbol{\alpha}_1 + k_2\boldsymbol{\alpha}_2 + k_3\boldsymbol{\alpha}_4$，其中 $k_1, k_2, k_3$ 为任意常数
\\
\noindent (C) $\boldsymbol{X} = k_1\boldsymbol{\alpha}_1 + k_2\boldsymbol{\alpha}_3 + k_3\boldsymbol{\alpha}_4$，其中 $k_1, k_2, k_3$ 为任意常数
\\
\noindent (D) $\boldsymbol{X} = k_1\boldsymbol{\alpha}_2 + k_2\boldsymbol{\alpha}_3 + k_3\boldsymbol{\alpha}_4$，其中 $k_1, k_2, k_3$ 为任意常数

\vfill
\newpage
\qn{2020.8}
设 $\boldsymbol{A}$ 为 3 阶矩阵，$\boldsymbol{\alpha}_1, \boldsymbol{\alpha}_2$ 为 $\boldsymbol{A}$ 的属于特征值 1 的线性无关的特征向量，$\boldsymbol{\alpha}_3$ 为 $\boldsymbol{A}$ 的属于特征值 $-1$ 的特征向量，则使得 $\boldsymbol{P}^{-1}\boldsymbol{A}\boldsymbol{P} = \begin{pmatrix} 1 & 0 & 0 \\ 0 & -1 & 0 \\ 0 & 0 & 1 \end{pmatrix}$ 的可逆矩阵 $\boldsymbol{P}$ 为（　　）．
\\
\noindent (A) $(\boldsymbol{\alpha}_1 + \boldsymbol{\alpha}_3, \boldsymbol{\alpha}_2, -\boldsymbol{\alpha}_3)$
\\
\noindent (B) $(\boldsymbol{\alpha}_1 + \boldsymbol{\alpha}_2, \boldsymbol{\alpha}_2, -\boldsymbol{\alpha}_3)$
\\
\noindent (C) $(\boldsymbol{\alpha}_1 + \boldsymbol{\alpha}_3, -\boldsymbol{\alpha}_3, \boldsymbol{\alpha}_2)$
\\
\noindent (D) $(\boldsymbol{\alpha}_1 + \boldsymbol{\alpha}_2, -\boldsymbol{\alpha}_3, \boldsymbol{\alpha}_2)$

\vfill
\newpage
\qn{2020.9}
设 $\begin{cases} x = \sqrt{t^2+1}, \\ y = \ln(t + \sqrt{t^2+1}), \end{cases}$ 则 $\left.\dfrac{\mathrm{d}^2 y}{\mathrm{d}x^2}\right|_{t=1} =$ ______．

\vfill
\newpage
\qn{2020.10}
$\displaystyle\int_{0}^{1}\mathrm{d}y\int_{\sqrt{y}}^{1} \sqrt{x^3+1}\,\mathrm{d}x =$ ______．

\vfill
\newpage
\qn{2020.11}
设 $z = \arctan[xy + \sin(x+y)]$，则 $\mathrm{d}z\big|_{(0,\pi)} =$ ______．

\vfill
\newpage
\qn{2020.12}
斜边长为 $2a$ 的等腰直角三角形平板铅直地沉入水中，且斜边与水面相齐，记重力加速度为 $g$，水密度为 $\rho$，则该平板一侧所受的水压力为 ______．

\vfill
\newpage
\qn{2020.13}
设 $y = y(x)$ 满足 $y'' + 2y' + y = 0$，且 $y(0) = 0$，$y'(0) = 1$，则 $\displaystyle\int_{0}^{+\infty} y(x)\,\mathrm{d}x =$ ______．

\vfill
\newpage
\qn{2020.14}
行列式 $\begin{vmatrix} a & 0 & -1 & 1 \\ 0 & a & 1 & -1 \\ -1 & 1 & a & 0 \\ 1 & -1 & 0 & a \end{vmatrix} =$ ______．

\vfill
\newpage
\qn{2020.15}
（本题满分 10 分）
求曲线 $y = \dfrac{x^{1+x}}{(1+x)^x}\ (x > 0)$ 的斜渐近线方程．

\vfill
\newpage
\qn{2020.16}
（本题满分 10 分）
已知函数 $f(x)$ 连续且 $\lim_{x\to 0} \dfrac{f(x)}{x} = 1$，$g(x) = \displaystyle\int_{0}^{1} f(xt)\,\mathrm{d}t$，求 $g'(x)$，并证明 $g'(x)$ 在 $x = 0$ 处连续．

\vfill
\newpage
\qn{2020.17}
（本题满分 10 分）
求函数 $f(x,y) = x^3 + 8y^3 - xy$ 的极值．

\vfill
\newpage
\qn{2020.18}
（本题满分 10 分）
设函数 $f(x)$ 的定义域为 $(0, +\infty)$ 且满足 $2f(x) + x^2 f\left(\dfrac{1}{x}\right) = \dfrac{x^2 + 2x}{\sqrt{1+x^2}}$，求 $f(x)$，并求曲线 $y = f(x)$，$y = \dfrac{1}{2}$，$y = \dfrac{\sqrt{3}}{2}$ 及 $y$ 轴所围图形绕 $x$ 轴旋转所成旋转体的体积．

\vfill
\newpage
\qn{2020.19}
（本题满分 10 分）
平面区域 $D$ 由直线 $x = 1$，$x = 2$，$y = x$ 与 $x$ 轴围成，计算 $\displaystyle\iint_D \dfrac{\sqrt{x^2+y^2}}{x}\,\mathrm{d}x\mathrm{d}y$．

\vfill
\newpage
\qn{2020.20}
（本题满分 11 分）
设 $f(x) = \displaystyle\int_{1}^{x} \mathrm{e}^{t^2}\,\mathrm{d}t$．
\\
（Ⅰ）证明：存在 $\xi \in (1,2)$，使得 $f(\xi) = (2-\xi)\mathrm{e}^{\xi^2}$；
\\
（Ⅱ）证明：存在 $\eta \in (1,2)$，使得 $f(2) = \ln 2 \cdot \eta\mathrm{e}^{\eta^2}$．

\vfill
\newpage
\qn{2020.21}
（本题满分 11 分）
设曲线 $y = f(x)$ 可导，且 $f'(x) > 0$，曲线 $y = f(x)\ (x \geqslant 0)$ 经过坐标原点 $O$，其上任意一点 $M$ 处的切线与 $x$ 轴交于 $T$，又 $MP$ 垂直 $x$ 轴于点 $P$，已知由曲线 $y = f(x)$，直线 $MP$ 以及 $x$ 轴所围图形的面积与 $\triangle MTP$ 的面积之比恒为 $3:2$，求满足上述条件的曲线方程．

\vfill
\newpage
\qn{2020.22}
（本题满分 11 分）
设二次型 $f(x_1, x_2, x_3) = x_1^2 + x_2^2 + x_3^2 + 2ax_1x_2 + 2ax_1x_3 + 2ax_2x_3$ 经可逆线性变换 $\begin{pmatrix} x_1 \\ x_2 \\ x_3 \end{pmatrix} = \boldsymbol{P}\begin{pmatrix} y_1 \\ y_2 \\ y_3 \end{pmatrix}$ 化为二次型 $g(y_1, y_2, y_3) = y_1^2 + y_2^2 + 4y_3^2 + 2y_1y_2$．
\\
（Ⅰ）求 $a$ 的值；
\\
（Ⅱ）求可逆矩阵 $\boldsymbol{P}$．

\vfill
\newpage
\qn{2020.23}
（本题满分 11 分）
设 $\boldsymbol{A}$ 为 2 阶矩阵，$\boldsymbol{P} = (\boldsymbol{\alpha}, \boldsymbol{A}\boldsymbol{\alpha})$，其中 $\boldsymbol{\alpha}$ 是非零向量且不是 $\boldsymbol{A}$ 的特征向量．
\\
（Ⅰ）证明：$\boldsymbol{P}$ 为可逆矩阵；
\\
（Ⅱ）若 $\boldsymbol{A}^2\boldsymbol{\alpha} + \boldsymbol{A}\boldsymbol{\alpha} - 6\boldsymbol{\alpha} = \boldsymbol{0}$，求 $\boldsymbol{P}^{-1}\boldsymbol{A}\boldsymbol{P}$，并判断 $\boldsymbol{A}$ 是否相似于对角矩阵．

\vfill
\newpage
\end{document}